%======================================================================
\section*{How this work was done}\addcontentsline{toc}{section}{How this work was done}
%======================================================================

\paragraph{Roles.} This paper is the work of its two authors together with AI models from more than one developer. The authors set the
direction: they chose the question, fixed the scope, decided which lines to pursue and which to drop, set the verification protocol, and
take full responsibility for every claim. The models did most of the exploration, derivation, computation and drafting. For the group and the game
they proposed constructions and proof strategies, wrote and checked proofs, and carried out the computations. They found the qubit count
of \cref{sec:count}, the transfer and return arguments of \cref{sec:transfer}, the rate for CHSH$(n)$ and the Wasserstein-$1$ argument for
the group $K$ of Slofstra and Vidick. They also wrote the verification programs, searched the literature, and drafted and edited the
text.

\paragraph{Sources.} The paper merges and supersedes two earlier unpublished drafts of ours. The first built the game and proved that every strategy
needs local dimension $\exp(\eps^{-1/2+o(1)})$, and every strategy on a maximally entangled state $\exp(\eps^{-1+o(1)})$; the second
developed the state-weighted method and closed that gap. This version proves every step from the definitions, apart from results of other
authors that are cited with their numbers. In particular it reproves the facts about Houghton's group that it uses (\cref{app:dehn}), and
it does not depend on our preprints~\cite{DM-lamps,DM-H3}, which appear only in remarks and credits. Material of the drafts that the main
results do not need has been left out: the general question of when Slofstra's embedding transfers approximate representations, a
comparison map between the solution groups of different games, the entropy needed by the Slofstra--Vidick game, and a discussion of the
three-player game of Ji, Leung and Vidick.

\paragraph{Verification.} The checks fall into three kinds.
\begin{itemize}[itemsep=1pt]
\item \emph{Machine checks of finite statements.} The programs described in \cref{app:cert} check the following:
\begin{itemize}[itemsep=0pt]
\item the presentation $\Pres$ and its model in $\GPp$;
\item the $29$ derivations behind \cref{prop:presentation}(d), with the area bound $A_*=60$;
\item the $27$ words of $\RP$;
\item the linear system of the game;
\item the area certificates behind $A_0=42$.
\end{itemize}
Each statement is checked by at least two programs that share no code, and every checking program runs negative
controls; \cref{app:cert} says which program establishes what. Their recorded output is in \texttt{anc/}.
\item \emph{Numerical attacks on the analytic statements.} A program in the ancillary files,
\texttt{anc/numerics/analytic\_checks.py}, evaluates both sides of the qubit count (\cref{thm:count}), the rounding and sequential-measurement lemmas, the loss bound of the return theorem (\cref{thm:return}) on
perturbed models of the magic square game, and the bound of \cref{thm:K-profile}, on random and optimized instances with flat, full-rank,
low-rank and pure states, and runs a negative control showing that the state commutators of~\ref{prelim:S2} cannot be dropped from the return
theorem. No inequality was violated. In the earlier drafts the triple construction for CHSH$(n)$ and the Wasserstein-$1$ contraction for
$K$ were attacked in the same way. A further program, \texttt{anc/sv\_balanced\_lift.py}, builds the lift of \cref{lem:sv-lift} for small
$m$, evaluates every relator, and confirms, as a negative control, that the literal lifts cannot be balanced.
\item \emph{Reading in the earlier drafts.} The qubit count, the transfer theorem, the argument for CHSH$(n)$ and a first form of the
argument for $K$ were proof-read by a model instance other than the one that derived them. They were then reviewed blind, refutation first,
by a model of a different developer that had not seen our working files. That review found no counterexample to any principal statement
and two repairable gaps, which are repaired here. The group-theoretic part was checked adversarially in the first draft.
\item \emph{Reading of this version.} Two model instances that had not taken part in writing this version read it in full, refutation
first: one read the analytic half and one the group-theoretic half. They re-derived the proofs, recomputed the constants, checked the cited
results of other papers against their sources, and attacked the analytic lemmas numerically. They also rebuilt the presentation and the
game independently of the shipped programs. They found no error in a main statement. They corrected the following slips, none of which
affects the main results:
\begin{itemize}[itemsep=0pt]
\item an intermediate defect bound in the proof of \cref{lem:gadget};
\item the range of $\eps$ in \cref{thm:sv-game} and in Theorem~E;
\item a remark claiming a case of equality;
\item the base of a logarithm in the introduction.
\end{itemize}
They also found that $17$ generators of an earlier version of $\Pres$ played no role. These have been removed, which shrinks the game from
$2061$ to $1687$ equations.
\item \emph{A blind audit.} A model of another developer, which had seen none of our working files, then audited this version refutation
first. It re-derived the main arguments, including the return direction (\cref{thm:return}), which no model of another developer had read
before; reran every program; and checked every reference against its primary source. It found no counterexample to a main statement. It
found that the proof of the upper bound of~\cite[Thm.~1.1]{SV18}, which Theorem~F cites, uses a lift from~\cite[Prop.~4.8]{Slo19} that,
for a group with two non-involutive generators such as $K$, does not leave the marked involution balanced as the last step of that proof
needs, and \cref{lem:sv-lift} now supplies a balanced lift for the explicit models of $K$; that our
account of the machine checks claimed more than the programs established, which led to the second cell-count verifier for $A_*=60$ and to
the cap of $42$ factors in both verifiers of (C4); and several smaller points. All are corrected here.
\end{itemize}
A formalization in Lean is in progress and is not part of this version. The proofs are given in full, and the checks support them but do not replace them.
